\section{Reference Implementation}
\label{sec:design}

To evaluate our analytical results, we implement \system{}, a
detection-and-mitigation framework that identifies faulty nodes and
reduces their routing weight. As shown in
Section~\ref{sec:complementarity}, the specific response strategy is
inconsequential under P2C---what matters is detection speed. We
describe \system{} primarily to document the detection mechanism used
in our evaluation.

\subsection{Architecture Overview}

\system{} operates as a control plane alongside the P2C data plane.
Every \emph{decision epoch} $\Delta$ (default 0.5\,s), it executes a
three-stage pipeline: \textbf{Monitor} $\to$ \textbf{Detect} $\to$
\textbf{Act}.

\paragraph{Monitor.}
Collects per-node metrics from the data plane: request departure
intervals, queue depths, and system-wide throughput. Computes
\emph{spare capacity} as the gap between observed throughput and
theoretical maximum, smoothed with an exponential moving average
($\alpha = 0.3$).

\paragraph{Detect.}
Estimates per-node service time from the P10 of inter-departure
intervals---a load-invariant signal derived from M/D/1 queueing theory
(Section~\ref{sec:detection}). Computes severity as:
\begin{equation}
  \sigma_i = \max\left(0,\ 1 - \frac{\tbase}{d_i}\right)
  \label{eq:severity}
\end{equation}
where $d_i$ is the estimated service time of node $i$. A 5$\times$
slowdown yields $\sigma = 0.8$; a 10$\times$ yields $\sigma = 0.9$.

\paragraph{Act.}
Selects a per-node strategy based on severity thresholds and applies
the corresponding traffic adjustment.

\subsection{Detection: Departure Interval Estimation}
\label{sec:detection}

A critical challenge in slow fault detection is distinguishing genuine
degradation from transient queueing effects. At high load ($\rho >
0.85$), all nodes exhibit elevated response times, causing latency-based
detectors to produce false positives.

We exploit a property of M/D/1 queues: when the queue is continuously
busy (which occurs most of the time at high utilization), inter-departure
times equal the actual service time, regardless of queue depth. The P10
percentile of departure intervals filters out idle gaps (when the queue
empties) while capturing the service-time floor.

This signal is \emph{load-invariant}: a node with $\tbase = 10$\,ms
produces $d \approx 10$\,ms whether the system is at 70\% or 95\% load.
A 5$\times$ slowdown produces $d \approx 50$\,ms at any load level.

To prevent transient noise from triggering mitigation, we require
\emph{persistence confirmation}: severity must exceed a minimum
threshold ($\sigma > 0.05$) for two consecutive epochs before a node is
flagged as faulty.

\subsection{Response Strategies}
\label{sec:strategies}

As shown in Section~\ref{sec:complementarity}, under P2C the specific
response strategy has negligible impact on P99 ($< 2$\,ms difference).
\system{} implements three strategies activated by severity thresholds
($\theta_{\text{spec}} < \theta_{\text{shed}} < \theta_{\text{iso}}$):
speculate (send hedged copy to a healthy backup for mildly degraded
nodes), shed (reduce routing weight
$w_i = \max(w_{\min},\, 1 - \sigma_i \cdot \text{spare\_factor})$
for moderately degraded nodes), and isolate (fully exclude severely
degraded nodes, subject to capacity guard $\rhoeff \leq U$).
Isolation is blocked when remaining healthy capacity would exceed $U$.

\subsection{Load-Aware Gating}

The key design principle is that \emph{every mitigation action has a
capacity cost}. Hedging doubles request load; shedding redistributes
load to healthy nodes; isolation removes capacity entirely. At high
system load, these costs can outweigh the benefit, causing the
mitigation trap.

\system{} addresses this through the spare capacity signal, which gates
all three strategies:
\begin{itemize}
\item Hedge budget $\propto$ spare capacity (zero hedges when spare $= 0$).
\item SHED weight reduction $\propto$ spare\_factor (no shedding when
  spare $< 0.08$).
\item ISOLATE blocked when remaining $\rhoeff > U$.
\end{itemize}

This ensures that \system{} gracefully degrades: at low load, it uses
aggressive isolation for fastest recovery; at high load, it uses
conservative shedding to preserve capacity; at overload, it backs off
entirely rather than making things worse.
